\chapter{Section Formulation \Paper{5/6}}\label{ch:section-theory}

\section{Introduction}


\paragraph{Literature}
% Exact rotations plus warping
\cite{simo1991geometricallyexact} introduced a geometrically exact formulation which incorporates one cross-sectional warping mode in a Type 1 formulation. 
This was investigated further by \cite{gruttmann2000theory} who emphasized the use of sections intersecting the reference curve at an arbitrary point. 
This theory has also been investigated by \cite{rong2020geometrically} in the context of isogeometric analysis who considered the nonlinear Wagner term. However, in this work the reference curve is restricted to intersect sections at the centroid. 

\section{Kinematics}

The extended displacement field for the newly introduced section material points $\hat{\rho}$  is given by:
\begin{equation}
\hat{\bm{x}}(\reflenx, \hat{\rho}) 
= \bm{x}(\reflenx) + \FrameRotation \bigl(\PlanePosition (\hat{\rho}) + (\bm{\alpha}\cdot\bm{\varphi})\, \FrameBasis\bigr)
\label{eq:section-displacement}
\end{equation}
%
where the 2D vector $\PlanePosition$ with coordinates $y\FrameBasis_{2} + z\FrameBasis_3$ identifies a section material point $\hat{\rho}$ with respect to the reference line of the beam. 
This generalizes the kinematic assumption of \cite{simo1991geometricallyexact} and aligns with the work by \cite{klinkel2003anisotropic} and \cite{elfatmi2007nonuniform}.

For shear-bending, in $(x-z)$-plane, for example, the warping function is often represented by an odd cubic function of $z$ (Reddy et al., 1997; Wang et al., 2000; Eisenberger, 2003), and the warping parameter $\alpha$ may be either independent (Martinez et al., 1999; Tsai and Kelly, 2005) or taken as the shear strain $\gamma$ (Dufort et al., 2001).

Note that the origin of this coordinate system is not assumed to coincide with special point in the section. In what follows, variants of the symbol \(\bm{\SectionLocation}\) are used to designate a fixed point in the cross section with respect to the section coordinate system. 
%


\subsection{Deformation}

For \(\hat{\FramePosition} = \hat{\chi}(\xi,\rho)\), consider the tangent map \(\bm{F} \triangleq T\hat{\chi}\) which is commonly known as the \emph{deformation gradient}.
% \subsection{Exact Kinematics}
%
% It will be convenient to introduce the relative displacement \(\hat{\bm{r}}\) given by the expression:
% $$
% \hat{\bm{r}} \triangleq \FrameRotation^{\mathrm{t}} (\hat{\bm{x}} - \bm{x}) = \PlanePosition  + (\bm{\alpha}\cdot\bm{\varphi})\,\FrameBasis
% $$
The \emph{exact} deformation gradient for the extended material field $\hat{\bm{x}}$ given by \cref{eq:section-displacement} is:
$$
\bm{F}(\xi, \hat{\rho}) 
% = \nabla \hat{\bm{x}} 
= \FrameRotation\left[\bm{1} + \bm{\varepsilon}\otimes\FrameBasis\right] + \FrameRotation\nabla \bm{\omega}
$$
where
$$
\bm{\varepsilon}(\hat{\rho}) \triangleq\ShearStrain  + \CurveStrain  \times (\PlanePosition  + \bm{\omega})
\qquad\text{ and }\qquad
\bm{\omega} \triangleq (\bm{\alpha}\cdot\bm{\varphi})\, \FrameBasis
$$
and $\ShearStrain $ and $\CurveStrain $ are the primary beam strain variables.
Note that in the presence of warping $\bm{\varepsilon}$ is nonlinear in
the strain variables $\FrameStrain$ and may be written as the sum of linear $\bar{\bm{\varepsilon}}$ and quadratic $\CurveStrain \times\bm{\omega}$ terms:
$$
\bm{\varepsilon} = \bar{\bm{\varepsilon}} + \CurveStrain \times\bm{\omega}
\qquad\text{ with }\qquad
\bar{\bm{\varepsilon}} \triangleq \ShearStrain  + \CurveStrain \times\PlanePosition 
$$
%
The exact Green-Lagrange strain tensor may be expressed as:
\begin{equation}
\SolidStrain = \bm{\varepsilon}
\stackrel{\mathrm{s}}{\otimes}
\FrameBasis 
+\nabla\omega \stackrel{\mathrm{s}}{\otimes} \FrameBasis
+({\bm{\varepsilon}}\cdot\FrameBasis)\,\nabla\omega \stackrel{\mathrm{s}}{\otimes} \FrameBasis
+ \frac{1}{2}\left(
  \bm{\varepsilon}\cdot\bm{\varepsilon}\,\FrameBasis\otimes\FrameBasis
+ \nabla\omega\otimes\nabla\omega
\right)
\label{eq:exact-gl-strain}
\end{equation}
where the symmetric bun product \(\bm{a}\stackrel{\mathrm{s}}{\otimes}\bm{b}\) is defined for vectors \(\bm{a}\), \(\bm{b}\), and \(\bm{c}\) by
\begin{equation*}
    (\bm{a}\stackrel{\mathrm{s}}{\otimes}\bm{b} )\bm{c} \triangleq \frac{1}{2} \bigl(\bm{a}\,(\bm{b}\cdot\bm{c}) + \bm{b}\,(\bm{a}\cdot\bm{c})\bigr)
\end{equation*}
%
In the absence of warping one has $\bm{\omega}=\bm{0}$ and $\bm{\varepsilon}\equiv \bar{\bm{\varepsilon}}$:
$$
\SolidStrain = \bar{\bm{\varepsilon}}
\stackrel{\mathrm{s}}{\otimes}
\FrameBasis 
+ \frac{1}{2}\left.
\bar{\bm{\varepsilon}}\cdot\bar{\bm{\varepsilon}}\,\FrameBasis\otimes\FrameBasis\right. \\
$$

\subsection{Truncation}
Even for geometrically exact models, it is common to assume linear section kinematics (see, e.g., \cite{simo1991geometricallyexact}). A common approach is to assume \(\bm{\varepsilon} \approx \bar{\bm{\varepsilon}}\) and 
\(\bm{\varepsilon}\cdot\bm{\varepsilon} \approx
  \tau^2 (\bm{r}\cdot\bm{r})\). 
This reduces \cref{eq:exact-gl-strain} to
%
\begin{equation}
\SolidStrain \approx \bar{\bm{\varepsilon}}
\stackrel{\mathrm{s}}{\otimes}
\FrameBasis 
+\nabla\omega \stackrel{\mathrm{s}}{\otimes} \FrameBasis
+ \tfrac{1}{2}\left.\varkappa^2\,
\bm{r}\cdot\bm{r}\,\FrameBasis\otimes\FrameBasis
\right.
\label{eq:wagner-gl-strain}
\end{equation}
which is similar to that used by \cite{hsiao2000corotational,battini2002corotational,jonker2021threedimensional} and \cite{rong2020geometrically}. 
Linearizing \cref{eq:wagner-gl-strain} furnishes
\[
\delta \bm{E} = \delta \ShearStrain - \bm{r}\times\delta \CurveStrain + (\FrameBasis\otimes \bm{\varphi})\delta \bm{\alpha} + (\nabla \bm{\varphi})^{\mathrm{t}} \delta \bm{\alpha}' + \varkappa r^2\, \FrameBasis\otimes\FrameBasis\delta \varkappa 
\]
where \(r^2 \triangleq \bm{r}\cdot\bm{r}\). Rearranging
\begin{equation*}
    \delta \bm{E} = \hat{\bm{B}}_{\mathrm{L}}\delta \bm{e} +  \hat{\bm{B}}_{\mathrm{W}}(\varkappa) \delta \bm{e}
\end{equation*}
with
\begin{equation}
\hat{\bm{B}}_{\mathrm{L}}  = 
\left[\bm{1} 
~~~ -\PlanePosition^{\times}
~~~ \FrameBasis\otimes\bm{\varphi}
~~~ \FrameBasis_{\beta} \otimes \nabla_{\beta} \bm{\varphi}
\right]
\quad
\text{and}
\quad
\hat{\bm{B}}_{\mathrm{W}} = \left[\bm{0} 
~~~ \varkappa r^2\,\FrameBasis\otimes\FrameBasis  
~~~ \bm{0}
~~~ \bm{0}
\right]
\end{equation}
%
\paragraph{Remark}
\(\hat{\bm{B}}_{\mathrm{L}}\) is equivalent to \cite[Eq. 2.59]{lecorvec2012nonlinear}, who uses the notation
\[
\bm{a}_{x}^w = \FrameBasis \otimes \bm{\varphi},
\qquad\text{ and }\qquad 
\bm{a}_{yz}^w = \FrameBasis_{\beta} \otimes \nabla_{\beta} \bm{\varphi}
\]


\section{Resultants}
Consider a 3D Constitutive law relating the second Piola-Kirchhoff stress tensor \(\bm{S}\) to Green's Lagrangian strain tensor \(\SolidStrain\).

Recall the Frobenius inner product for two second order tensors \(\bm{A}\) and \(\bm{B}\) with coordinate representations \(A^{ij} \triangleq \mathbf{G}^i\cdot \bm{A}\mathbf{G}^j\) and  \(B_{ij} \triangleq \mathbf{G}_i\cdot \bm{B}\mathbf{G}_j\) on a basis \(\{\mathbf{G}_k\}\)
\begin{equation*}
    \left<\bm{A},\,\bm{B}\right> = \operatorname{tr} \bm{A}\bm{B}^{\mathrm{t}} = A^{ij} B_{ij}
\end{equation*}
If \(\delta \SolidStrain = \left(\hat{\bm{B}} \delta\FrameStrain\right) \stackrel{s}{\otimes}\FrameBasis\) for some second-order tensor \(\hat{\bm{B}}\), then
\begin{equation*}
    \left<\bm{S},\,\delta\SolidStrain\right> = \SectionIntegral{ 
    \operatorname{tr} \left[\bm{S}\frac{1}{2}\left(\hat{\bm{B}}\delta \FrameStrain \otimes \FrameBasis + \FrameBasis\otimes \hat{\bm{B}}\delta \FrameStrain \right)
    \right]
    }
\end{equation*}
Exploiting linearity of the trace together with the identity \(\operatorname{tr} (\bm{a}\otimes\bm{b}) = \bm{a}\cdot \bm{b}\) for arbitrary vectors \(\bm{a}\) and \(\bm{b}\) it follows
\begin{equation*}
    \left<\bm{S},\,\delta\SolidStrain\right> = \delta \FrameStrain\cdot \SectionIntegral{
        \hat{\bm{B}}^{\mathrm{t}}\frac{1}{2}\left(\bm{S}\FrameBasis + \bm{S}^{\mathrm{t}}\FrameBasis\right)
    }
\end{equation*}
and consequently by requiring \(\left<\FrameResult,\,\delta\FrameStrain\right>=\left<\bm{S},\,\delta\SolidStrain\right>\)
\begin{equation*}
\FrameResult \equiv \SectionIntegral{\hat{\bm{B}}^{\mathrm{t}}\frac{1}{2}\left(\bm{S}\FrameBasis + \bm{S}^{\mathrm{t}}\FrameBasis\right)}
\end{equation*}
For symmetric \(\bm{S}\) one further has
\begin{equation}
\FrameResult = \SectionIntegral{
  \hat{\bm{B}}^{\mathrm{t}} \bm{S}\FrameBasis
}
\end{equation}
This has the intuitive implication that ony the first ``column'' and/or ``row'' of the stress tensor \(\bm{S}\) should contribute to the frame stress resultants \(\ShearResult\) and \(\CurveResult\).

The stress tensor can be expressed as follows
\begin{equation}
\bm{S}=\sigma \FrameBasis \otimes \FrameBasis+\bm{\tau} \otimes \FrameBasis+\FrameBasis \otimes \bm{\tau}
\label{eq:saint-venant-stress}
\end{equation}
\paragraph{Linearization}
\begin{equation*}
    D {\FrameResult} [\delta \FrameStrain] = \hat{\bm{B}}^{\mathrm{t}}D \bm{S} [D \SolidStrain [\delta \FrameStrain]]
\end{equation*}
Using \(\mathbb{C}[\delta \SolidStrain] \triangleq D\bm{S}[\delta \SolidStrain]\) for vectors \(\delta \SolidStrain\) tangent at \(\SolidStrain\) 
\begin{equation*}
 D {\FrameResult} [\delta \FrameStrain] = \hat{\bm{B}}^{\mathrm{t}} \mathbb{C} D \SolidStrain [\delta \FrameStrain]
\end{equation*}
Now using \(D \SolidStrain [\delta \FrameStrain] = (\hat{\bm{B}}\delta \FrameStrain) \stackrel{s}{\otimes} \FrameBasis\)
\begin{equation}
\hat{\mathbf{K}} = \SectionIntegral{
    \hat{\mathbf{B}}^{\mathrm{t}} \mathbb{C}\hat{\bm{B}} + \hat{\bm{G}}
    }
\label{eq:section-tangent}
\end{equation}

\subsection{Truncation}
When the truncated strain \cref{eq:wagner-gl-strain} is used, it is convenient to rewrite \cref{eq:section-tangent} as the sum of explicit \textbf{Linear} \(\hat{\bm{K}}_{\mathrm{L}}\) and \textbf{Wagner} \(\hat{\bm{K}}_{\mathrm{W}}\) terms given by
\begin{equation}
\hat{\bm{K}}_{\mathrm{L}} = \SectionIntegral{
\left[\begin{matrix}
\mathbb{C} & - \mathbb{C} \bm{r}^{\times} & \mathbb{C} \FrameBasis\otimes \bm{\varphi} & \mathbb{C} \FrameBasis_{\beta} \otimes \nabla_{\beta} \bm{\varphi}\\
%\bm{r}^{\times} \mathbb{C} 
& - \bm{r}^{\times} \mathbb{C} \bm{r}^{\times} 
& \bm{r}^{\times} \mathbb{C} \FrameBasis\otimes \bm{\varphi} & \bm{r}^{\times} \mathbb{C} \FrameBasis_{\beta} \otimes \nabla_{\beta} \bm{\varphi}\\
% \bm{\varphi} \otimes \FrameBasis \mathbb{C} 
& % - \bm{\varphi} \otimes \FrameBasis \mathbb{C} \bm{r}^{\times} 
& \bm{\varphi} \otimes \FrameBasis \mathbb{C} \FrameBasis\otimes \bm{\varphi} 
& \bm{\varphi} \otimes \FrameBasis \mathbb{C} \FrameBasis_{\beta} \otimes \nabla_{\beta} \bm{\varphi}\\
\mathrm{sym.} % \nabla_{\beta} \bm{\varphi} \otimes \FrameBasis_{\beta} \mathbb{C} 
& % - \nabla_{\beta} \bm{\varphi} \otimes \FrameBasis_{\beta} \mathbb{C} \bm{r}^{\times} 
& % \nabla_{\beta} \bm{\varphi} \otimes \FrameBasis_{\beta} \mathbb{C} \FrameBasis\otimes \bm{\varphi} 
& \nabla_{\beta} \bm{\varphi} \otimes \FrameBasis_{\beta} \mathbb{C} \FrameBasis_{\beta} \otimes \nabla_{\beta} \bm{\varphi}
\end{matrix}\right]
}
\label{eq:section-linear-tangent}
\end{equation}
and using \(r^2 \triangleq \bm{r}\cdot\bm{r}\)
\begin{equation}
\begin{aligned}
\hat{\bm{K}}_{\mathrm{W}} =
\kappa_T & \SectionIntegral{
\left[\begin{matrix}
\bm{0} & r^{2}  \mathbb{C} \FrameBasis \otimes \FrameBasis & \bm{0} & \bm{0} \\
r^{2}  \FrameBasis \otimes \FrameBasis \mathbb{C} 
& r^{2}  \left(\bm{r}^{\times} \mathbb{C} \FrameBasis \otimes \FrameBasis -  \FrameBasis \otimes \FrameBasis \mathbb{C} \bm{r}^{\times}\right) & r^{2}  \FrameBasis \otimes \FrameBasis \mathbb{C} \FrameBasis\otimes \bm{\varphi} & r^{2}  \FrameBasis \otimes \FrameBasis \mathbb{C} \FrameBasis_{\beta} \otimes \nabla_{\beta} \bm{\varphi}\\
\bm{0} & r^{2}  \bm{\varphi} \otimes \FrameBasis \mathbb{C} \FrameBasis \otimes \FrameBasis & \bm{0} & \bm{0} \\
\bm{0} & r^{2}  \nabla_{\beta} \bm{\varphi} \otimes \FrameBasis_{\beta} \mathbb{C} \FrameBasis \otimes \FrameBasis & \bm{0} & \bm{0} 
\end{matrix}\right]} \\
+\kappa_T^2 & \SectionIntegral{\left[\begin{matrix}
\bm{0} & \bm{0} & \bm{0} & \bm{0} \\
\bm{0} & r^{4} \FrameBasis \otimes \FrameBasis \mathbb{C} \FrameBasis \otimes \FrameBasis & \bm{0} & \bm{0} \\
\bm{0} & \bm{0} & \bm{0} & \bm{0} \\
\bm{0} & \bm{0} & \bm{0} & \bm{0}
\end{matrix}\right]
} \\
&+\SectionIntegral{\left[\begin{matrix}
\bm{0} & \bm{0} & \bm{0} & \bm{0} \\
\bm{0} & r^2 \sigma & \bm{0} & \bm{0} \\
\bm{0} & \bm{0} & \bm{0} & \bm{0} \\
\bm{0} & \bm{0} & \bm{0} & \bm{0} 
\end{matrix}\right]
}
\end{aligned}
\label{eq:section-wagner-tangent}
\end{equation}

\subsection{Change of Reference}
\begin{equation*}
    \bm{u}_A(\bm{r}_A) = \bm{u}_A(\bm{0}) + \bm{\theta}\times\bm{r}_A + \bm{\alpha}\cdot\bm{\varphi}_A(\bm{r}_A) \,\FrameBasis
\end{equation*}
%
\begin{equation*}
\bm{u}_A(\bm{0}) = \bm{u}_B(\bm{0}) + \bm{\theta}\times \bm{\SectionLocation_{AB}}
\end{equation*}
%
Thus
\begin{equation*}
\begin{pmatrix}
\delta\bm{\gamma}_A \\ \delta \bm{\kappa}_A
\end{pmatrix}
= \begin{bmatrix}
    \bm{1} & -\bm{\SectionLocation_{AB}^{\times}} \\
    \bm{0} & \bm{1}
\end{bmatrix}
\begin{pmatrix}
\delta \bm{\gamma}_B \\ \delta \bm{\kappa}_B
\end{pmatrix}
\end{equation*}


\subsection{Summary}

\vspace{1ex}
\begin{ambox}[Section]{box:section-type1}
Let a section be specified with point-wise data \(\mathcal{A}, \mathcal{M}, \bm{\varphi}\), and \(\nabla \bm{\varphi}\), and a selection of either \textbf{Linear}, \textbf{Wagner}, or \textbf{Finite} kinematics. 
\begin{itemize}
    \item Kinematics: At each point \(\hat{\rho}\) the strain \(\bm{E}\) is evaluated from \cref{eq:exact-gl-strain}, \cref{eq:wagner-gl-strain}, or \cref{eq:} and \(\hat{\bm{B}}\) from \cref{eq:}, \cref{eq:} or \cref{eq:}

    \item Constraints: With \(\mathbb{P}\) from \cref{eq:warping-projection}, augment \(\hat{\bm{B}}\)

    \item Response: \(\bm{s}\)
    \[
    \FrameResult= \SectionIntegral{\hat{\bm{B}}^{\mathrm{t}} \bm{S}\FrameBasis}
    \]

    \item Tangent: Integrate \(\hat{\bm{K}}\) using \cref{eq:section-tangent}
    \begin{equation*}
    \hat{\bm{K}} = \SectionIntegral{
        \hat{\bm{B}}^{\mathrm{t}} \mathbb{C}\hat{\bm{B}} + \hat{\bm{G}}}
    \end{equation*}
    where \(\hat{\bm{G}}\) vanishes under \textbf{Linear} kinematics, and is given by \cref{eq:section-wagner-geometric} for \textbf{Wagner} kinematics.

\end{itemize}
\end{ambox}

\section{Example: Isotropic Elasticity}
Consider the St. Venant-Kirchhoff material model with material parameters \(\lambda\) and \(\mu\)
\begin{equation*}
    \bm{S} = \lambda  \bm{1} \operatorname{tr} \bm{E} + 2 \mu \bm{E}
\end{equation*}
The fourth-order tangent tensors are
\begin{equation}
\FourTensor{C}=\kappa 3\FourTensor{P}_{\mathrm{vol}} + 2 \mu \FourTensor{P}_{\mathrm{dev}}
\qquad\text{ and }\qquad
\FourTensor{C}^{-1} = \left.\frac{1}{3\kappa} \FourTensor{P}_{\mathrm{vol}} + \frac{1}{2 \mu}\FourTensor{P}_{\mathrm{dev}}\right.
\label{eq:isotropic-moduli}
\end{equation}
where $\kappa=\lambda+\frac{2}{3} \mu>0$ is the bulk modulus, $\mu>0$ is the shear modulus, and
\begin{equation*}
\FourTensor{P}_{\mathrm{vol}}=\frac{1}{3} \bm{1} \otimes \bm{1}
,\quad 
\FourTensor{P}_{\mathrm{dev}}=\FourTensor{P}_{\mathrm{sym}}-\FourTensor{P}_{\mathrm{vol}} 
\quad \text { and } \quad 
\FourTensor{P}_{\mathrm{sym}} = \frac{1}{2}\FrameBasis_{i} \otimes \FrameBasis_{j} \otimes\left( \FrameBasis_{i} \otimes \FrameBasis_{j} + \FrameBasis_{j} \otimes \FrameBasis_{i}\right)
\end{equation*}
%
Using \cref{eq:isotropic-moduli} in \cref{eq:section-linear-tangent} reveals 
the tensors \(\LinearMM\), \(\LinearWW\), \(\LinearVV\), \(\SectionNW\), \(\SectionNV\), \(\LinearMW\) and \(\LinearMV\) which encapsulate the isotropic response of a homogeneous \textbf{Linear} section 
\begin{equation}
\left[
\begin{matrix}
    A \bm{1} 
  & A\bar{\bm{\pi}}^{\times} % nm
  & \SectionNW  % E
  & \SectionNV  % G
  \\

%
  & \LinearMM
  & \LinearMW % E
  & \LinearMV % G
  \\

%
  & %
  & \LinearWW % E?
  & \\

%
  & %
  & %
  & \LinearVV
\end{matrix}
\right]
= \SectionIntegral{
\hat{\bm{B}}^{\mathrm{t}}_{\mathrm{L}} \hat{\bm{B}}_{\mathrm{L}} 
}
= \SectionIntegral{
\left[
\begin{matrix}
\bm{1} 
  &-\bm{r}^{\times} 
  & \FrameBasis\otimes \bm{\varphi} 
  & \FrameBasis_{\beta} \otimes \nabla_{\beta} \bm{\varphi} \\

%- \bm{r}^{\times} 
  & - \left[\bm{r}^{\times}\right]^{2} 
  & \bm{r} \times \FrameBasis \otimes \bm{\varphi} 
  & \bm{r} \times \FrameBasis_{\beta} \otimes \nabla_{\beta} \bm{\varphi} \\

%\bm{\varphi} \otimes \FrameBasis
  & %[\bm{\varphi} \otimes \FrameBasis] \bm{r}^{\times} 
  & \bm{\varphi} \otimes \bm{\varphi}
  & \bm{0} \\

%
  & % 
  & %
  & \nabla_{\beta} \bm{\varphi} \otimes \nabla_{\beta} \bm{\varphi}
\end{matrix}
\right]
}
\label{eq:def-linear-isotropic-coupling-tensors}
\end{equation}
Similarly, for the higher-order \textbf{Wagner} kinematics one additionally has from \cref{eq:isotropic-moduli} in \cref{eq:section-wagner-tangent} the constants
\begin{equation}
\begin{bmatrix}
\bm{0} & \tau \bm{A}_{r^2} & \bm{0} & \bm{0} 
\\
& \tau \bm{I}_{r^2} + \tau^2 \bm{I}_{r^4}
& \tau \bm{R}_{r^2} 
& \bm{0} \\
& & \bm{0} & \bm{0} \\
& & & \bm{0}
\end{bmatrix}
\triangleq \
\SectionIntegral{\left[\begin{matrix}
\bm{0} & \tau r^{2} \FrameBasis \otimes \FrameBasis & \bm{0} & \bm{0}\\
\tau r^{2} \FrameBasis \otimes \FrameBasis 
& \tau r^{2} (\bm{r} {\times} \FrameBasis \otimes \FrameBasis +  \FrameBasis \otimes \bm{r} {\times}\FrameBasis) + \tau^2 r^{4} \FrameBasis \otimes \FrameBasis & \tau r^{2} \FrameBasis\otimes \bm{\varphi} & \bm{0} \\
\bm{0} & \tau r^{2}  \bm{\varphi} \otimes \FrameBasis  & \bm{0} & \bm{0} \\
 &  &  &  \bm{0}
\end{matrix}
\right]}
\end{equation}
with 
\(\bm{A}_{r^2} \triangleq A_{r^2} \FrameBasis \otimes \FrameBasis\),
\(\bm{I}_{r^4} \triangleq I_{r^4} \FrameBasis \otimes \FrameBasis\).

% \subsubsection{Linear Kinematics}

\paragraph{$\bm{n}$-$\bm{n}$}

\paragraph{m-m}\label{m-m}

\[
\begin{aligned}
\LinearMM &= \SectionIntegral{
     [\bm{r}^\times][\bm{r}^\times]^{\mathrm{t}} 
}
= \SectionIntegral{ \left[\begin{matrix}
y^{2} + z^{2} & 0 & 0\\
0 & z^{2} & - z y \\
0 & - z y & \,\,y^{2}
\end{matrix}\right]
}
% &= \int (\bm{r}\cdot\bm{r})\bm{1} - \bm{r}\otimes \bm{r}\, \mathrm{d}A \\
% &=\left[\delta_{\alpha \beta} \mathbf{1}-\FrameBasis_\alpha \otimes \FrameBasis_\beta\right]\int_{\Omega} \xi_\alpha \xi_\beta \, \mathrm{d}A  \\
\\
% &\triangleq I_0 \FrameBasis \otimes \FrameBasis + I_{\alpha \beta} \FrameBasis_\alpha \otimes \FrameBasis_\beta  \\
\end{aligned}
\]
\textbf{three} independent constants.
% \[
% \begin{aligned}
% I_0    &\triangleq  y^2 + z^2 \\
% I_y    &\triangleq  I_{22} = \int z^2\, \mathrm{d} A \\
% I_z    &\triangleq  I_{33} = \int y^2\, \mathrm{d} A \\
% I_{zy} &\triangleq -I_{23} = \int yz \, \mathrm{d} A
% \end{aligned}
% \]
Note
\[
\LinearMM = \SectionIntegral{\bm{r}\cdot\bm{r}\bm{1} - \bm{r}\otimes\bm{r}}
\]

\begin{itemize}
    \item The \emph{principal axes} of the cross section are defined as the coordinate system which diagonalizes \(\LinearMM\). Consequently, in this coordinate system the resultant moments \(M_{\beta}\) are often decoupled. 
\end{itemize}

\paragraph{m-m Wagner}

\[
\tau r^{2} (\bm{r} {\times} \FrameBasis \otimes \FrameBasis +  \FrameBasis \otimes \bm{r} {\times}\FrameBasis)
\]

\paragraph{w-w}\label{w-w}

\begin{equation}
\LinearWW = \SectionIntegral{\bm{\varphi}\otimes\bm{\varphi}}
= \SectionIntegral{ \begin{bmatrix}
\varphi^2  & \varphi \psi_y & \varphi \psi_z \\
           & \psi_y^2 & \psi_z \psi_y \\
           &          & \psi_z^2
\end{bmatrix}}
= \begin{bmatrix}
J_{w} & 0 & 0 \\
0 & C_{\psi} & 0 \\
0 & 0 & C_{\psi}
\end{bmatrix}
\end{equation}
where \textbf{6} independent constants are defined. This is equivalent to the \emph{warping matrix} of \cite{elfatmi2007nonuniform}.

\begin{itemize}
\tightlist
\item
  If \(\bm{\varphi} = \tilde{\varphi} \,\FrameBasis\) then \(J_w\) is the warping
  constant of \cite[p39]{vlasov1961thinwalled} 
  \[
  J_w = \SectionIntegral{ \tilde{\varphi}^2 }
  \]
  Some other names for this constant include \emph{sectorial moment of inertia} \cref{vlasov1961thinwalled}, 
\end{itemize}

\paragraph{\(\bm{v}-\bm{v}\)}

\begin{equation}
\LinearVV = \SectionIntegral{\nabla_{\beta} \bm{\varphi}\otimes\nabla_{\beta} \bm{\varphi}}
= \SectionIntegral{ 
\begin{bmatrix}
\varphi_{,y}^2 & \varphi_{,y}\psi_{y,y} & \varphi_{,y}\psi_{z,y} \\
& \psi_{y,y}^2 & \psi_{y,y}\psi_{z,y} \\
& & \psi_{z,y}^2
\end{bmatrix} +
\begin{bmatrix}
\varphi_{,z}^2 & \varphi_{,z}\psi_{y,z} & \varphi_{,z}\psi_{z,z} \\
& \psi_{y,z}^2 & \psi_{y,z}\psi_{z,z} \\
& & \psi_{z,z}^2
\end{bmatrix}
}
\end{equation}
\textcite{gruttmann2000theory} refer to \(J_{v}\) 
% NOTE: I previously used C_{\alpha} for J_v
with the notation ``\(I_0 + I_T\)''
\[
  J_{v}  = \SectionIntegral{ \tilde{\varphi}_{, y}^2 + \tilde{\varphi}_{, z}^2} = J_S - J
\]

\paragraph{Coupling in $\bm{n}-\bm{m}$}\label{n-m}

\[
\LinearNM
\triangleq \SectionIntegral{ -[\bm{r}^\times] }
= \SectionIntegral{ \begin{bmatrix}
 0 & z &-y \\
-z & 0 & 0 \\
 y & 0 & 0
\end{bmatrix} \, \mathrm{d} A
 = \begin{bmatrix}
     0 & Q_{y} & -Q_{z} \\
-Q_{y} & 0 & 0 \\
 Q_{z} & 0 & 0
\end{bmatrix}
}
\]
where \textbf{two} independent constants known as the \emph{first moments of area} have been defined.

\begin{itemize}
\item
  These are exactly an area-scaling of the centroid coordinates \(\CentroidLocation\) defined in \cref{eq:define-centroid}, that is,  
  \[
  \bm{A}_r \equiv A \CentroidLocation^{\times}
  \]
\item
  If the reference curve passes through the centroid as
  in \cite{simo1991geometricallyexact} then \(\bm{A}_r = \bm{0}\)
\end{itemize}

\paragraph{Coupling in $\bm{n}-\bm{w}$}\label{n-w}

\[
\SectionNW \triangleq \SectionIntegral{ \FrameBasis \otimes\bm{\varphi} 
}
= \SectionIntegral{\begin{bmatrix}
\varphi & \psi_y & \psi_z \\
0 & 0 & 0 \\
0 & 0 & 0 \\
\end{bmatrix}
}
% = \begin{bmatrix}
% J_{nw} & S_{ny} & S_{nz} \\
% 0 & 0 & 0 \\
% 0 & 0 & 0
% \end{bmatrix}
\] 
where \textbf{three} independent constants have been defined and
\[
\SectionNV \triangleq \SectionIntegral{ \FrameBasis_{\beta}\otimes \nabla_{\beta}\bm{\varphi}
}
= \SectionIntegral{\begin{bmatrix}
0 & 0 & 0 \\
\varphi_{,y} & \psi_{y,y} & \psi_{z,y} \\
\varphi_{,z} & \psi_{y,z} & \psi_{z,z} \\
\end{bmatrix}
}
= \begin{bmatrix}
0 &     0 & 0 \\
J_{ny} & S_{nyy} & S_{nzy} \\
J_{nz} & S_{nyz} & S_{nzz}
\end{bmatrix}
\] 
assuming \(\bm{\varphi}' = \bm{0}\).

\begin{itemize}
\item If \(\varphi\) is defined by a pure Neumann problem of the Laplacian, then the condition 
\[
\SectionIntegral{\varphi} = 0
\]
is most likely satisfied by \(\varphi\). This is generally the case for St. Venant warping.
\item
  If \(\varphi\) is referred to the \emph{shear
  center} (e.g., \cite{gruttmann1998geometrical}) 
  \begin{equation*}
  J_{ny} = \int \tilde{\varphi}_{,y} = -(\tilde{z} - \CentroidLocationZ) A
  \qquad\text{and}\qquad
  J_{nz} = \int \tilde{\varphi}_{,z} =  (\tilde{y} - \CentroidLocationY) A
  \end{equation*}
\item
  If \(\varphi\) is referred to the \emph{centroid}
  (i.e.~\(\varphi = \hat{\varphi}\) like \cite{gruttmann2000theory}) then:
  \[
  J_{ny} = \int \hat{\varphi}_{,y} =    \CentroidLocationZ A \\
  \qquad\text{and}\qquad
  J_{nz} = \int \hat{\varphi}_{,z} =  - \CentroidLocationY A
  \]
  where \(\CentroidLocationZ\) and \(\CentroidLocationY\) are the coordinates of the centroid.
\end{itemize}

\paragraph{m-w}\label{m-w}
% NOTE: May remove the symbols J_my, etc... if they are always equal to -Iv, then they are redundant.
\[
\LinearMW 
= \SectionIntegral{ \bm{r} \times\FrameBasis \otimes \bm{\varphi} } 
= \SectionIntegral{ \begin{bmatrix}
0 & 0 & 0 \\
~z \varphi & ~z \psi_{y} & ~z \psi_{z} \\
-y \varphi & -y \psi_{y} & -y \psi_z
\end{bmatrix}
}
= \begin{bmatrix}
0 & 0 & 0 \\
J_{my} & S_{myy} & S_{mzy} \\
J_{mz} & S_{myz} & S_{mzz}
\end{bmatrix}
\]
introducing \textbf{two} shear-interacting constants \(J_{my}\) and \(J_{mz}\) for torsion.
\begin{itemize}
\tightlist
\item
  If \(\varphi\) is referred to the centroid 
  \[
  \begin{aligned}
  S_y &= \int z \hat{\varphi} \, \mathrm{d} A =   \hat{I}_{z} \hat{\SectionLocation}_z - \hat{I}_{yz} \hat{y} \\
  S_z &= \int y \hat{\varphi} \, \mathrm{d} A = - \hat{I}_{y} \hat{\SectionLocation}_y + \hat{I}_{yz} \hat{z}
  \end{aligned}
  \]
\end{itemize}
For moment-bishear interaction:
\[
\LinearMV
= \SectionIntegral{ \bm{r} \times\FrameBasis_{\beta} \otimes \nabla_{\beta}\bm{\varphi} } 
= \SectionIntegral{ \begin{bmatrix}
-z \varphi_{, y} + y \varphi_{, z} &
-z \psi_{y, y} + y \psi_{y, z} &
-z \psi_{y, y} + y \psi_{y, z} \\
0 & 0 & 0 \\
0 & 0 & 0 \\
\end{bmatrix}
}
% = \begin{bmatrix}
% J_{mw} & S_{my} & S_{mz} \\
% 0 & 0 & 0 \\
% 0 & 0 & 0 \\
% \end{bmatrix}
\]
Generally, however, \(\LinearMV \equiv - \LinearVV\).
% \[
% \SectionMW 
% = \SectionIntegral{ \bm{r} \times\FrameBasis \otimes \bm{\varphi} } 
% = \SectionIntegral{ \begin{bmatrix}
% 0 & -z \varphi_{, y} + y \varphi_{, z} \\
%  z \varphi & 0 \\
% -y \varphi & 0
% \end{bmatrix}
% }
% = \begin{bmatrix}
% 0 & S_{\alpha} \\
% S_{y} & 0 \\
% S_{z} & 0
% \end{bmatrix}
% \]




The final linear stiffness matrix for a general isotropic cross section is given explicitly by \cref{eq:ks-isotropic}.

\begin{landscape}
\topskip0pt
\vspace*{\fill}
\begin{equation}
\hat{\bm{K}} = 
\left[\begin{array}{ccc|ccc|ccc|ccc}
EA & 0 & 0 & 0 & E Q_z & - E Q_y & E Q_{\varphi} & E Q_{y} & E Q_{z} & 0 & 0 & 0\\
0 & GA & 0 & - G Q_z & 0 & 0 & 0 & 0 & 0 & G J_{n y} & G Q_{y y} & G Q_{z y} \\
0 & 0 & GA & G Q_y & 0 & 0 & 0 & 0 & 0 & G J_{n z} & G Q_{y z} & G Q_{z z} \\ \hline
0 & - G Q_z & G Q_y & G I_0 & 0 & 0 & 0 & 0 & 0 & G R_{\varphi} &  G R_y & G R_z \\
E z & 0 & 0 & 0 & E I_{z z} & - E I_{y z} & E J_{my} & E \psi_{y} z & E \psi_{z} z & 0 & 0 & 0\\
- E y & 0 & 0 & 0 & - E I_{y z} & E I_{y y} & E J_{m z} & - E \psi_{y} y & - E \psi_{z} y & 0 & 0 & 0\\ \hline
E \varphi & 0 & 0 & 0 & E \varphi z & - E \varphi y & E \varphi^{2} & E \psi_{y} \varphi & E \psi_{z} \varphi & 0 & 0 & 0\\
E \psi_{y} & 0 & 0 & 0 & E \psi_{y} z & - E \psi_{y} y & E \psi_{y} \varphi & E \psi_{y}^{2} & E \psi_{y} \psi_{z} & 0 & 0 & 0\\
E \psi_{z} & 0 & 0 & 0 & E \psi_{z} z & - E \psi_{z} y & E \psi_{z} \varphi & E \psi_{y} \psi_{z} & E \psi_{z}^{2} & 0 & 0 & 0\\ \hline
0 & G \varphi_{,y} & G \varphi_{,z} & G . & 0 & 0 & 0 & 0 & 0 & G S_{\varphi \varphi} & G S_{\varphi y} & G S_{\varphi z} \\
0 & G \psi_{y,y} & G \psi_{y,z} & G . & 0 & 0 & 0 & 0 & 0 & G S_{y \varphi} & G S_{y y} & G S_{y z} \\
0 & G \psi_{z,y} & G \psi_{z,z} & G . & 0 & 0 & 0 & 0 & 0 & G S_{z \varphi} & G S_{zy} & G S_{zz}
\end{array}\right]
\label{eq:ks-isotropic}
\end{equation}
\vspace*{\fill}
\end{landscape}
% \end{minipage}
% \end{sideways}

\section{Summary and Remarks}

\begin{enumerate}
\item \textbf{Torsion Warping}: In this case only a single warping mode is considered so that \(\bm{\alpha}\) is a scalar and the warping mode \(\varphi\) is St. Venant's warping function.
\iffalse
$$
\nabla \omega = \bm{\Phi}\begin{pmatrix}\alpha^{\prime} \\ \alpha \end{pmatrix}
\qquad\text{ with }\qquad
\bm{\Phi} \triangleq \begin{bmatrix}
\varphi & 0 \\ 0 & \varphi_{, y} \\ 0 & \varphi_{, z}
\end{bmatrix}
$$
\fi
In the absence of constraints on \(\alpha\), this corresponds to the RBV formulation described by \cite{armero2021modeling}.
This model is well suited for implementation in \gls{frame:shear} models. Examples of finite element implementations include \cite{simo1991geometricallyexact} and \cite{gruttmann2000theory}.

$$
  \mathbf{L}_{\mathrm{n}}
= \mathbf{L}_{\mathrm{m}}
= \mathbf{L}_{\mathrm{w}}
= \bm{0}
$$
$$
\hat{\bm{B}}\mathbf{P} = \hat{\bm{B}} = \left[\begin{array}{ccc|ccc|ccc|ccc}
1 & 0 & 0 & 0 & z & - y & \varphi & \psi_{y} & \psi_{z} & 0 & 0 & 0 \\
0 & 1 & 0 & - z & 0 & 0 & 0 & 0 & 0 & \varphi_{,y} & \psi_{y,y} & \psi_{z,y}\\
0 & 0 & 1 & y & 0 & 0 & 0 & 0 & 0 & \varphi_{,z} & \psi_{y,z} & \psi_{z,z}
\end{array}\right]
$$
See Equation~(26) of \cite{dire2019mixed}

\item For Type 1 torsion and uniform (Type 3) shear warping \cref{eq:ks-isotropic} reduces to:
\begin{equation*}
\hat{\mathbf{K}} = \left[\begin{array}{ccc|ccc|cc}
E A & 0 & 0 & 0 & E Q_{y} & E Q_{z} & E R_{w} & 0\\
0 & G A_{y} & 0 & G Q_{yx} & 0 & 0 & 0 & G R_{y}\\
0 & 0 & G A_{z} & G Q_{zx} & 0 & 0 & 0 & G R_{z}\\ \hline
0 & G Q_{yx} & G Q_{zx} & G I_{0} & 0 & 0 & 0 & G S_{a} \\
E Q_{y} & 0 & 0 & 0 & E I_{y} & - E I_{yz} & E S_{y} & 0 \\
E Q_{z} & 0 & 0 & 0 & - E I_{yz} & E I_{z} & E S_{z} & 0 \\ \hline
E R_{w} & 0 & 0 & 0 & E S_{y} & E S_{z} & E C_{w} & 0 \\
0 & G R_{y} & G R_{z} & G S_{a} & 0 & 0 & 0 & G C_{a}
\end{array}\right]
\end{equation*}

\item T2 (constrained) torsion. Referred to as the TWKV formulation by \cite{armero2021modeling}, this formulation imposes the constraint 
  \[
  \alpha = \CurveStrain \cdot \FrameBasis
  \]
  This assumption derives from observing that, for very thin elements, uniform warping displacements can be considered across the thickness of the cross-section membratures and, thus, transverse shear strains in this direction can be neglected [20,21]. This is usually adequate for open cross-sections, but can lead to incorrect solutions for closed profiles.
  Examples include \cite{nukala2004mixed} and \cite{zhang2011formulation} for symmetric sections, whose work is implemented in OpenSees and derives from \cite{alemdar2001distributed}.

\item T2 (constrained) torsion and T2 shear, linearized

$$
\mathbf{L}_{\mathrm{n}} = \FrameBasis_{\beta} \otimes \FrameBasis_{\beta}, \quad
\mathbf{L}_{\mathrm{m}} = \FrameBasis_{} \otimes \FrameBasis_{},
\quad\text{ and }\quad
\mathbf{L}_{\mathrm{w}} = \bm{1}
$$
% $$
% \mathbf{L}_{\mathrm{p}} = \left[\begin{matrix}
% 0 & 0 & 0 & 1 & 0 & 0\\
% 0 & 1 & 0 & 0 & 0 & 0\\
% 0 & 0 & 1 & 0 & 0 & 0
% \end{matrix}\right]
% \qquad\text{ and }\qquad
% \mathbf{L}_{\mathrm{w}} = \bm{1}
% $$
$$
\hat{\bm{B}}\mathbf{P} = \left[\begin{matrix}
1 & 0 & 0 & 0 & z & - y & \varphi & \psi_{y} & \psi_{z}\\
0 & \psi_{y,y} + 1 & \psi_{z,y} & \varphi_{,y} - z & 0 & 0 & 0 & 0 & 0\\
0 & \psi_{y,z} & \psi_{z,z} + 1 & \varphi_{,z} + y & 0 & 0 & 0 & 0 & 0\end{matrix}\right]
$$
This is quivalent to Equation~(29) for the ``EVDE'' formulation of  \cite{addessi2021enriched}. 
\item T2 (constrained) torsion and T3 (uniform) shear, linearized

$$
\mathbf{L}_{\mathrm{n}} = \FrameBasis_{\beta} \otimes \FrameBasis_{\beta}, \quad
\mathbf{L}_{\mathrm{m}} = \FrameBasis_{} \otimes \FrameBasis_{},
\quad\text{ and }\quad
\mathbf{L}_{\mathrm{w}} = \bm{1}
$$

$$
\hat{\bm{B}}\mathbf{P} = \left[\begin{matrix}
1 & 0 & 0 & 0 & z & - y & \varphi \\
0 & \psi_{y,y} + 1 & \psi_{z,y} & \varphi_{,y} - z & 0 & 0 & 0 \\
0 & \psi_{y,z} & \psi_{z,z} + 1 & \varphi_{,z} + y & 0 & 0 & 0 
\end{matrix}\right]
$$
Described by \cite[Section 3.1.2]{lecorvec2012nonlinear} with \(k_y \beta  \triangleq 1+\psi_{y,y}\). This model was also employed by \cite[Equation 3.51]{dire20173d}.
For a homogeneous section with linear-elastic St.Venant-Kirchhoff material  and \textbf{Linear} section kinematics one has:
\begin{equation*}
\hat{\mathbf{K}} = \left[\begin{array}{ccc|ccc|c}
E A & 0 & 0 & 0 & E Q_{y} & E Q_{z} & E R_{w}\\
0 & G A_{y} & 0 & G (Q_{yx} + R_{y}) & 0 & 0 & 0\\
0 & 0 & G A_{z} & G (Q_{zx} + R_{z}) & 0 & 0 & 0\\ \hline
0 & G Q_{yx} + G R_{y} & G Q_{zx} + G R_{z} & G (C_{a} + I_{0} + 2 S_{a}) & 0 & 0 & 0\\
E Q_{y} & 0 & 0 & 0 & E I_{y} & - E I_{yz} & E S_{y}\\
E Q_{z} & 0 & 0 & 0 & - E I_{yz} & E I_{z} & E S_{z}\\ \hline
E R_{w} & 0 & 0 & 0 & E S_{y} & E S_{z} & E C_{w}
\end{array}\right]
\end{equation*}
%
Note that this generalizes the work of \cite{saritas2006mixed}, as presented in Equation~3.23 of \cite{dire20173d}.

\item T3 \emph{Uniform torsion} In this case warping is assumed to depend on the torsional component of the curvature $\CurveStrain $ so that $\alpha = \FrameBasis\cdot\CurveStrain $ and $\alpha^{\prime}=0$. Therefore:
$$
\bm{\omega} = (\varphi \FrameBasis\otimes\FrameBasis)\CurveStrain  = \tau \varphi \, \FrameBasis
$$
and
$$
\bm{\varepsilon} = \ShearStrain  + \CurveStrain \times\PlanePosition  + \tau\varphi \CurveStrain \times\FrameBasis
$$

\item Type 3 (uniform) torsion and Type 3 shear warping
$$
\mathbf{L}_{\mathrm{n}} = \FrameBasis_{\beta} \otimes \FrameBasis_{\beta}, \quad
\mathbf{L}_{\mathrm{m}} = \FrameBasis_{} \otimes \FrameBasis_{},
\quad\text{ and }\quad
\mathbf{L}_{\mathrm{w}} = \bm{1}
$$
$$
\hat{\bm{B}}\mathbf{P} = \left[\begin{matrix}
1 & 0 & 0 & 0 & z & - y \\
0 & \psi_{y,y} + 1 & \psi_{z,y} & \varphi_{,y} - z & 0 & 0 \\
0 & \psi_{y,z} & \psi_{z,z} + 1 & \varphi_{,z} + y & 0 & 0
\end{matrix}\right]
$$
%
See also \citep{wackerfuss2009mixed}
\begin{equation*}
\left[\begin{array}{cccccc}
E A & 0 & 0 & 0 & E A s_3 & -E A s_2 \\
0 & \kappa_2 G A & 0 & -\kappa_2 G A m_3 & 0 & 0 \\
0 & 0 & \kappa_3 G A & \kappa_3 G A m_2 & 0 & 0 \\
0 & -\kappa_2 G A m_3 & \kappa_3 G A m_2 & \widetilde{I}_T & 0 & 0 \\
E A s_3 & 0 & 0 & 0 & E \widetilde{I}_{33} & -E \widetilde{I}_{32} \\
-E A s_2 & 0 & 0 & 0 & -E \widetilde{I}_{23} & E \widetilde{I}_{22}
\end{array}\right]
\end{equation*}

\item Type 3 torsion and Type 4 shear  \citep{gruttmann1998geometrical} where additionally \(\psi_{y,y}  \approx \psi_{z,z} \approx 0\)
$$
\hat{\bm{B}}\mathbf{P} = \left[\begin{matrix}
1 & 0 & 0 & 0 & z & - y \\
0 & 1 & 0 & \varphi_{,y} - z & 0 & 0 \\
0 & 0 & 1 & \varphi_{,z} + y & 0 & 0
\end{matrix}\right]
$$

\item Type 4 torsion and Type 3 shear. Described by \cite{lecorvec2012nonlinear} as the \emph{Shear beam} model.
$$
\hat{\bm{B}}\mathbf{P} = \left[\begin{matrix}
1 & 0 & 0 & 0 & z & - y \\
0 & 1+\psi_{y,y} & 0 & - z & 0 & 0 \\
0 & 0 & 1+\psi_{z,z} &   y & 0 & 0
\end{matrix}\right]
$$
\item \textbf{Plane Section}: In the absence of warping one has $\bm{\omega}=\bm{0}$ and $\bm{\varepsilon}\equiv \bar{\bm{\varepsilon}}$:
$$
\SolidStrain = \bar{\bm{\varepsilon}}
\stackrel{\mathrm{s}}{\otimes}
\FrameBasis 
+ \frac{1}{2}\left.
\bar{\bm{\varepsilon}}\cdot\bar{\bm{\varepsilon}}\,\FrameBasis\otimes\FrameBasis\right. \\
$$


\end{enumerate}


\newpage
\section{A: Constraints}
\newcommand{\ThreeTensor}[1]{\mathsf{#1}}

The ramification of the rigid cross section assumption is that the normal and shear strain components in the plane of a cross section vanish; in our coordinate system that means
\(E_{yy}=E_{zz}=0\) and \(E_{yz}=E_{zy}=0\). 
For isotropy, this
constraint is not a problem in the shear strain components because the
shear stresses and strains are uncoupled under isotropy (for example,
\(S_{yz}=2 \mu E_{yz}\) ). Therefore, vanishing shear strains simply
implies vanishing shear stress.

The constitutive equations for isotropic hyperelasticity have an
inherent coupling of the normal strains and stresses. The stresses are
given by \[
\mathbf{S}=\lambda \operatorname{tr}(\bm{E}) \bm{1}+2 \mu \mathbf{E}
\] Since under the strain hypothesis
\(\operatorname{tr}(\bm{E})=E_{xx}\), we have
\(S_{yy}=S_{zz}=\lambda E_{xx}\). Now, from the constitutive law
\(S_{xx}=(\lambda+2 \mu) E_{xx}\) so \(S_{yy}=S_{zz}=\nu S_{xx}\), where
\(\nu\) is Poisson's ratio. Thus, the constraint induces normal stresses
in the plane of the cross section owing to Poisson's effect (tension in
one direction causes lateral contraction of the dimensions perpendicular
to the direction of tension). Observational evidence on the behavior of
beams would indicate that these stresses tend to be rather small. In
fact, it is possible that the presence of these stresses will violate
the traction boundary conditions on the lateral surface of the beam. For
a beam with a traction-free lateral surface, we can argue that the
normal stresses \(S_{yy}\) and \(S_{zz}\) are very small because the
cross-sectional dimensions are small compared with the length of the
beam. We would like to make the assumption that \(S_{yy}=S_{zz}=0\).

Let us examine what would have happened if we had made the assumption of
vanishing normal stress and not the assumption of vanishing normal
strain. We made exactly that assumption for the uniaxial tension test in
order to recast the constitutive equations in terms of Young's modulus
and Poisson's ratio. When we made the assumption that
\(S_{yy}=S_{zz}=0\), we got the constitutive relationship
\(S_{xx}=C E_{xx}\), rather than \(S_{xx}=(\lambda+2 \mu) E_{xx}\),
which results for the vanishing strain assumption. This gives us a way
to partially recover from our difficulties. In the constitutive
relationships for beams, the quantity \(E=\lambda+2 \mu\) appears
repeatedly. If we simply substitute the value of Young's modulus \(E=C\)
instead, then the results of beam theory accord well with observation.

\subsection{Stress Space}\label{stress-space}

Consider \(\mathbb{S}_{\mathrm{rod}}\subset \mathbb{S}\) 
\[
\operatorname{dim}\mathbb{S}_{\mathrm{rod}} = 6 - 3 = 3
\]

\subsubsection{Coordinates}\label{coordinates}

Let \(\bm{\sigma} \triangleq \bm{S}\FrameBasis\) and define
\(\bm{\varepsilon}\) so that
\(\FourTensor{P}_{\mathrm{rod}}\bm{E} = \FourTensor{P}_{\mathrm{sym}}[\FrameBasis \otimes\bm{\varepsilon}]\).
Note that the shear components of \(\bm{\varepsilon}\) and \(\bm{E}\)
are related by 
\[
\FrameBasis_{\beta}\cdot\bm{\varepsilon} = \FrameBasis_{\beta} \cdot 2\bm{E}\FrameBasis
\] 
which reads, in matrix notation,
\((\varepsilon_2, \varepsilon_3) = (2E_{21},\, 2E_{31})\).

The components \(\sigma_k\) of \(\bm{\sigma}\) on the frame basis
\(\{\FrameBasis_k\}\) are equivalent to the components of
\(\bm{S}\in\mathbb{S}_{\mathrm{rod}}\) on a basis \(\{\bm{I}_k\}\)
defined by \[
\bm{I}_k \triangleq \FrameBasis\otimes\FrameBasis_k
\]
%
An alternative basis for \(\mathbb{S}_{\mathrm{rod}}\) is furnished by
the three tensors 
\[
\hat{\bm{I}} = \FrameBasis \otimes \FrameBasis,\, 
\quad\text{ and }\quad 
\hat{\bm{I}}_{\beta} = \tfrac{1}{2}\left(\FrameBasis \otimes\FrameBasis_{\beta} + \FrameBasis_{\beta}\otimes\FrameBasis\right)
\]

\subsubsection{Projections}\label{projections}

The \emph{frame stress} condition is defined by three constraints 
\[
\FrameBasis_{\alpha}\cdot\bm{S}\FrameBasis_{\beta} = 0
\]
\[
\FourTensor{P}_{\mathrm{sec}} = \FrameBasis_{\alpha}\otimes\FrameBasis_{\beta}\otimes\FrameBasis_{\alpha}\otimes\FrameBasis_{\beta}
\quad\text{ and }\quad
\FourTensor{P}_{\mathrm{rod}} = \mathbb{I} - \FourTensor{P}_{\mathrm{sec}}
\]
where 
\[
\mathbb{I} = \FrameBasis_j\otimes\FrameBasis_k \otimes \FrameBasis_j \otimes \FrameBasis_k 
\]
also let 
\[
\FourTensor{P}_{\mathrm{rod}} = \FourTensor{P}_{\mathrm{eul}} + \FourTensor{P}_{\mathrm{cos}}
\]
where 
\[
\FourTensor{P}_{\mathrm{eul}}\triangleq \FrameBasis\otimes\FrameBasis\otimes\FrameBasis\otimes\FrameBasis
\quad\text{ and }\quad 
\FourTensor{P}_{\mathrm{cos}} \triangleq \mathbb{I} - \FourTensor{P}_{\mathrm{eul}}
\]

\subsubsection{Moduli Components}\label{moduli-components}

Introduce a third order tensor
\(\ThreeTensor{P}_{\mathrm{sym}}^{\FrameBasis}\) defined so that 
\[
\ThreeTensor{P}_{\mathrm{sym}}^{\FrameBasis} \bm{\varepsilon} = \FourTensor{P}_{\mathrm{sym}}[\FrameBasis \otimes\bm{\varepsilon}] = \FourTensor{P}_{\mathrm{rod}}\bm{E}
\] 
We can find a representation of
\(\ThreeTensor{P}_{\mathrm{sym}}^{\FrameBasis}\) by appropriately
contracting \(\FourTensor{P}_{\mathrm{sym}}\) with \(\FrameBasis\). 
\[
\ThreeTensor{P}_{\mathrm{sym}}^{\FrameBasis} = \frac{1}{2} \left(\FrameBasis\otimes \FrameBasis_k\otimes\FrameBasis_k + \FrameBasis_k \otimes \FrameBasis\otimes\FrameBasis_k\right)
\] 
we should verify
\(\FourTensor{P}_{\mathrm{sec}}:\ThreeTensor{P}_{\mathrm{sym}}^{\FrameBasis} = \ThreeTensor{0}\).
Then we can write 
\[
\bm{S}\FrameBasis = \ThreeTensor{C}_{\mathrm{rod}}^{\FrameBasis}:\ThreeTensor{P}_{\mathrm{sym}}^{\FrameBasis} \bm{\varepsilon}
\]

\newpage
\section{St. Venant Torsional Warping}\label{sec:torsion-warping}

The objective of this section is to detail the computation of St. Venant's warping function \(\varphi\) and related cross section parameters \(C_{\varphi}, C_w, J, ...\).

\subsection{The Neumann Problem}
Consider a body with reference configuration 
\begin{equation*}
    \bm{x}_0(\reflenx) + \SectionCoordinate(\varsigma)
\end{equation*}
undergoing a deformation with the form
\begin{equation*}
    \bm{x}(\reflenx, \varsigma ) = \bm{x}_0(\reflenx) + \SectionCoordinate(\varsigma) + \vartheta \, \FrameBasis\times(\SectionCoordinate - \bm{\SectionLocation}) + \alpha \varphi(\varsigma) \FrameBasis
\end{equation*}
Subject to the traction-free boundary condition
\begin{equation*}
\begin{aligned}
    \bm{S}\mathbf{n} &= \bm{0}
\end{aligned}
\end{equation*}
where \(\mathbf{n}\) is the unit normal on the boundary in the section's plane. 
The displacement field \(\bm{u}\) is 
\begin{equation*}
\bm{u}(\reflenx, y, z) =
\vartheta \, \FrameBasis\times(\SectionCoordinate - \bm{\SectionLocation}) + \alpha \varphi(\SectionCoordinate) \FrameBasis
\end{equation*}
and the infinitesimal strain vector
\begin{equation*}
\bm{\varepsilon} = \vartheta' \FrameBasis \times (\SectionCoordinate - \bm{\SectionLocation}) + \alpha' \varphi \FrameBasis + \alpha \nabla \varphi
\end{equation*}
The traction-free boundary condition implies
\begin{equation*}
    \nabla_{\mathbf{n}}\bm{\varepsilon} = \bm{0}
\end{equation*}
Recall \cref{eq:saint-venant-stress} and that for a prismatic section the outward normal is orthogonal to the axial basis \(\FrameBasis\), \(\mathbf{n}\cdot \FrameBasis = 0\).
\begin{marsdenbox}
\noindent
Torsional warping problem for rotation about some point \(\bm{\SectionLocation}\)
\begin{equation}
\begin{array}{rll}
\Delta \varphi_{\SectionLocation} &=0 & \text { in } \Omega, \\
\nabla_{\mathbf{n}} \varphi_{\SectionLocation}  &=-\alpha \FrameBasis \times (\mathbf{r}-\bm{\SectionLocation}) \cdot \mathbf{n} & \text { on } \partial \Omega  \\
\SectionIntegral{\varphi_{\SectionLocation}} &= 0
\end{array}
\label{eq:torsion-neumann}
\end{equation}
where $\mathbf{n}$ is the unit outward normal to the boundary of the cross section domain.
\end{marsdenbox}


\subsection{Center of Twist}
A natural choice for \(\bm{\SectionLocation}\) is to select it such that the following conditions hold:
\begin{equation}
    \SectionIntegral{\varphi_{\SectionLocation} (\bm{r}-\bm{\SectionLocation})\times\FrameBasis} = \bm{0}
    \label{eq:shear-center-constraint}
\end{equation}
Thus, the point \(\tilde{\bm{\SectionLocation}}\) is defined as that for which the corresponding solution \(\tilde{\varphi}\) to \cref{eq:torsion-neumann} with \(\pi = \tilde{\bm{\SectionLocation}}\) satisfies \cref{eq:shear-center-constraint}. 
Note that \cref{eq:shear-center-constraint} simply requires that the \(m-w\) coupling tensor of \cref{eq:def-linear-isotropic-coupling-tensors} vanishes.
...
say
\[
\tilde{\varphi} \triangleq \varphi - [\tilde{\bm{\SectionLocation}} , \FrameBasis, (\bm{r} - \hat{\bm{\SectionLocation}})]
\]
In components 
\[
\tilde{\varphi} = \varphi - 
\]
or equivalently
\[
\tilde{\varphi}\bm{r}\times\FrameBasis = \varphi\bm{r}\times\FrameBasis - \bm{r}\times\FrameBasis \, \left(\tilde{\bm{\SectionLocation}} \cdot \FrameBasis \times \hat{\bm{r}}\right)
\]
\[
\tilde{\varphi}\bm{r}\times\FrameBasis = \varphi\bm{r}\times\FrameBasis - \bm{r}\times\FrameBasis \otimes \left(\FrameBasis \times \hat{\bm{r}}\right) \tilde{\bm{\SectionLocation}}
\]
Using this in \cref{eq:shear-center-constraint} furnishes the linear system:
\begin{equation}
    \SectionIntegral{\tilde{\varphi}\,\bm{r}\times\FrameBasis}
    =\left(\SectionIntegral{\bm{r}\times\FrameBasis\otimes\FrameBasis\times\bar{\bm{r}}} \right)\tilde{\bm{\SectionLocation}}
    \label{eq:derive-shear-center-1}
\end{equation}
where \(\bar{\bm{r}} \triangleq \bm{r} - \bar{\bm{\SectionLocation}}\).
Also recall that 
\[
\SectionIntegral{\varphi_{\SectionLocation} \bm{r}\times\FrameBasis} 
\equiv \SectionIntegral{\varphi_{\SectionLocation} (\bm{r}-\bar{\bm{\SectionLocation}})\times\FrameBasis} 
%\equiv \bm{0}
\]
Thus we use \(\bar{\bm{r}}\) for \(\bm{r}\) in \cref{eq:derive-shear-center-1}
\begin{equation}
    \SectionIntegral{\tilde{\varphi}\,\bar{\bm{r}}\times\FrameBasis}
    =\left(\SectionIntegral{\bar{\bm{r}}\times\FrameBasis\otimes\FrameBasis\times\bar{\bm{r}}} \right)\tilde{\bm{\SectionLocation}}
\end{equation}
Pulling out the constant \(\FrameBasis\) 
\begin{equation}
    \SectionIntegral{\tilde{\varphi}\,\bar{\bm{r}}\times\FrameBasis}
    =\FrameBasis^{\times}\left(\SectionIntegral{\bar{\bm{r}}\otimes\bar{\bm{r}}} \right)\FrameBasis^\times\tilde{\bm{\SectionLocation}}
\end{equation}


\subsection{Finite Elements}\label{finite-elements}

\cite{surana1979isoparametric}

\begin{quote}
\[
\begin{gathered}
\bm{K}_{a b}=\int \left(\frac{\partial N_{i}}{\partial x} \frac{\partial N_{b}}{\partial x}+\frac{\partial N_{i}}{\partial y} \frac{\partial N_{b}}{\partial y}\right) \operatorname{det}[J] \mathrm{d} \xi \mathrm{d} \eta \\
\bm{f}_{a}=\SectionIntegral{\left(y \frac{\partial N_{i}}{\partial x}-x \frac{\partial N_{i}}{\partial y}\right) \operatorname{det}[J]
} \\
\left\{\delta_{a}\right\}=\left\{\psi_{a}\right\}
\end{gathered}
\]
\end{quote}


\begin{equation}
    \bm{f}\cdot\bm{w} = \FrameBasis \cdot \SectionIntegral{(\bm{B}\bm{w} \times \FrameBasis_{\beta} \bm{N}\bm{x}^e_{\beta})}
\end{equation}

\begin{equation}
    \bm{f}\cdot\bm{w} = -\FrameBasis \cdot \FrameBasis_{\beta}\times \SectionIntegral{ (\bm{N}\bm{x}^e_{\beta})\, \bm{B}\bm{w}}
\end{equation}
\newpage
\section{A: Condensation}

\begin{equation}
\dot{\bm{S}} = \FourTensor{C}\dot{\bm{E}}
\label{eq:differential-constitution}
\end{equation}
We are concerned with constraints of the form
\begin{equation}
\FourTensor{P}_j \dot{\bm{S}} = \bm{0}
\label{eq:linear-stress-constraint}
\end{equation}
where \(\FourTensor{P}_j\) is a fourth-order projection tensor, ie, it has the property \(\FourTensor{P}_j^2 \equiv\FourTensor{P}_j\). 
Consider the decomposition
\begin{equation}
\bm{E} = \FourTensor{P}_i \bm{E} + \FourTensor{P}_j \bm{E}
\label{eq:strain-projection-decomposition}
\end{equation}
so that the nonzero stress rate \(\FourTensor{P}_i\dot{\bm{S}}\) is expressed as 
\begin{equation}
\FourTensor{P}_i\dot{\bm{S}}=\FourTensor{P}_i\FourTensor{C}\FourTensor{P}_i \bm{E} + \FourTensor{P}_i\FourTensor{C} \FourTensor{P}_j \bm{E}
\label{eq:nonzero-stress}
\end{equation}
Differentiating the constraint \cref{eq:linear-stress-constraint} furnishes with \cref{eq:differential-constitution}
$$
\bm{0} = \FourTensor{P}_j\FourTensor{C}\dot{\bm{E}} = \FourTensor{P}_j\FourTensor{C}\FourTensor{P}_i \bm{E} + \FourTensor{P}_j\FourTensor{C} \FourTensor{P}_j \bm{E}
$$
Moving the term with $\FourTensor{P}_j\bm{E}$ to the left produces
$$
\FourTensor{P}_j\FourTensor{C} \FourTensor{P}_j \FourTensor{P}_j \dot{\bm{E}} = -\FourTensor{P}_j \FourTensor{C}\FourTensor{P}_i \dot{\bm{E}}
$$
where $\FourTensor{P}_j$ is duplicated by exploiting the projection property $\FourTensor{P}_j^2 = \FourTensor{P}_j$. Solving the linear system furnishes
$$
\FourTensor{P}_j \dot{\bm{E}} = -(\FourTensor{P}_j\FourTensor{C}\FourTensor{P}_j)^{-1} \FourTensor{P}_j\FourTensor{C} \FourTensor{P}_i \dot{\bm{E}}
$$
Inserting this back in \cref{eq:strain-projection-decomposition} 
$$
\dot{\bm{E}} = \FourTensor{P}_i \dot{\bm{E}} - \FourTensor{P}_j(\FourTensor{P}_j\FourTensor{C}\FourTensor{P}_j)^{-1} \FourTensor{P}_j\FourTensor{C} \FourTensor{P}_i \dot{\bm{E}}
$$
Thus the unknown $\FourTensor{P}_j \dot{\bm{E}}$ can be eliminated from \cref{eq:nonzero-stress}
$$
\dot{\bm{S}} = \FourTensor{P}_i \dot{\bm{S}} = \left(\FourTensor{P}_i\FourTensor{C}\FourTensor{P}_i
-\FourTensor{P}_i\FourTensor{C}\FourTensor{P}_j(\FourTensor{P}_j\FourTensor{C}\FourTensor{P}_j)^{-1} \FourTensor{P}_j\FourTensor{C}\FourTensor{P}_i\right) \dot{\bm{E}}
$$

\newpage
\begin{equation}
\begin{aligned}
    \mathbf{I}_{\mathrm{eul}} & \triangleq \FrameBasis \otimes \FrameBasis \\
    \mathbf{I}_{\mathrm{cos}} & \triangleq \bm{1}-\mathbf{I}_{\mathrm{eul}} \\
\end{aligned}
\end{equation}

\[
\FourTensor{P}_{\mathrm{sec}} = \FrameBasis_{\alpha}\otimes\FrameBasis_{\beta}\otimes\FrameBasis_{\alpha}\otimes\FrameBasis_{\beta}
\quad\text{ and }\quad
\FourTensor{P}_{\mathrm{rod}} = \mathbb{I} - \FourTensor{P}_{\mathrm{sec}}
\]

$$
\begin{aligned}
\FourTensor{P}_{\mathrm{eul}}
&= \FrameBasis\otimes\FrameBasis\otimes\FrameBasis\otimes\FrameBasis = \mathbf{I}_{\mathrm{eul}}\otimes \mathbf{I}_{\mathrm{eul}}
\\
\FourTensor{P}_{\mathrm{vs}}
&= \frac{1}{3}\bm{1}\otimes\FrameBasis_{\alpha}\otimes\FrameBasis_{\alpha}
 = \frac{1}{3}\FourTensor{P}_{\mathrm{eul}} + \frac{1}{3}()
\\
\FourTensor{P}_{\mathrm{sv}}
&= \frac{1}{3}\FrameBasis_{\alpha}\otimes\FrameBasis_{\alpha}  \otimes \bm{1}
\\
\FourTensor{P}_{\mathrm{svs}} &
= \frac{1}{3}\FrameBasis_{\gamma}\otimes\FrameBasis_{\gamma}\otimes\FrameBasis_{\delta}\otimes\FrameBasis_{\delta}
\\
\FourTensor{P}_{\mathrm{svsv}}
&= \frac{2}{3} \FourTensor{P}_{\mathrm{sv}}
\\
\FourTensor{P}_{\mathrm{vsvs}} 
&=\frac{2}{3} \FourTensor{P}_{\mathrm{vs}}
%= \frac{2}{9}\FrameBasis_{k}\otimes\FrameBasis_{k}\otimes\FrameBasis_{\delta}\otimes\FrameBasis_{\delta} 
\\
\FourTensor{P}_{\mathrm{vsv}} 
    &= \frac{2}{9}\FourTensor{P}_{\mathrm{vol}}
\\
\FourTensor{P}_{\mathrm{vsvsv}} 
    &= \frac{4}{9}\FourTensor{P}_{\mathrm{vol}}
\\
\end{aligned}
$$
\begin{equation}
    \begin{aligned}
    \FourTensor{P}_{\mathrm{ds}} 
        &= \FourTensor{P}_{\mathrm{sec}} - \FourTensor{P}_{\mathrm{vs}} \\
    \FourTensor{P}_{\mathrm{dsv}} 
        &= \FourTensor{P}_{\mathrm{sv}} - \frac{2}{9}\FourTensor{P}_{\mathrm{vol}} \\
    \FourTensor{P}_{\mathrm{vsdsv}} 
        &= \frac{7}{9}\FourTensor{P}_{\mathrm{vol}} \\
    \FourTensor{P}_{\mathrm{dsd}} 
        &= \FourTensor{P}_{\mathrm{sec}} - \FourTensor{P}_{\mathrm{sv}} - \FourTensor{P}_{\mathrm{vs}} + \frac{2}{9}\FourTensor{P}_{\mathrm{vol}} \\
    \FourTensor{P}_{\mathrm{vsdsd}} 
        &=\frac{1}{3} \FourTensor{P}_{\mathrm{vs}} 
        - \frac{14}{81} \FourTensor{P}_{\mathrm{vol}}
        \\
    \FourTensor{P}_{\mathrm{dsds}} 
        &= \FourTensor{P}_{\mathrm{sec}} - \FourTensor{P}_{\mathrm{svs}} - \frac{7}{9}\FourTensor{P}_{\mathrm{vs}} \\
    \FourTensor{P}_{\mathrm{dsvs}} 
        &= \FourTensor{P}_{\mathrm{svs}} - \frac{2}{3}\FourTensor{P}_{\mathrm{vs}}  \\
    \FourTensor{P}_{\mathrm{dsvsv}}
        &= \frac{2}{3} \FourTensor{P}_{\mathrm{sv}} 
         - \frac{4}{27} \FourTensor{P}_{\mathrm{vol}} \\
    \FourTensor{P}_{\mathrm{dsvsd}}
        &= \FourTensor{P}_{\mathrm{svs}} 
         - \frac{2}{3} \FourTensor{P}_{\mathrm{vs}} 
         - \frac{2}{3} \FourTensor{P}_{\mathrm{sv}} 
         + \frac{4}{27} \FourTensor{P}_{\mathrm{vol}} \\
    \FourTensor{P}_{\mathrm{dsdsd}}
        &= \FourTensor{P}_{\mathrm{sec}} 
         - \frac{5}{3}\FourTensor{P}_{\mathrm{sv}}
         - \frac{7}{9} \FourTensor{P}_{\mathrm{vs}}
         - \FourTensor{P}_{\mathrm{svs}}
         + \frac{14}{81} \FourTensor{P}_{\mathrm{vol}}\\
    \FourTensor{P}_{\mathrm{dsdsv}}
        &= \frac{1}{3}\FourTensor{P}_{\mathrm{sv}}
         - \frac{14}{81} \FourTensor{P}_{\mathrm{vol}}
    \end{aligned}
\end{equation}